% GENERATED FILE. Edit canonical lessons, then run npm run generate:books.
% canonical-source-hash: fnv1a64:c01afcd3953a1a92
% canonical-lessons: SA-C280-sthanakam-kutra, SA-R280-asking-the-way

\chapter{Which Way to the Station?}
\label{ch:sa-which-way-to-the-station}
\hlchaptermodality{280}

\begin{chapteropening}
\textbf{By the end of this chapter:} \emph{I can ask the way (sthānakaṃ kutra asti? kathaṃ gacchāmi? kiyat dūram?) and follow the answer: left, right, straight on (vāmataḥ, dakṣiṇataḥ, ṛju gacchatu).}

Complete SA-R280-asking-the-way: ask the way and follow the answer.
\end{chapteropening}

\section[\texorpdfstring{\sk{स्थानकं} \sk{कुत्र} \sk{अस्ति}? (sthānakaṃ kutra asti?)}{स्थानकं कुत्र अस्ति? (sthānakaṃ kutra asti?)}]{\sk{स्थानकं} \sk{कुत्र} \sk{अस्ति}? (sthānakaṃ kutra asti?) — asking the way}

\label{lesson:SA-C280-sthanakam-kutra}



\begin{quote}
Before the new one: say \textquotedblleft{}where?\textquotedblright{} and \textquotedblleft{}please\textquotedblright{}. (\emph{kutra}, \emph{kṛpayā}.) You are in a town you do not know, and you need the station.
\end{quote}

\subsection*{You'll want to know: where, how, how far}
\begin{quote}
\textbf{\sk{स्थानकं} \sk{कुत्र} \sk{अस्ति}?}
\end{quote}

\begin{quote}
\emph{sthānakaṃ kutra asti?} — \textquotedblleft{}Where is the station?\textquotedblright{}
\end{quote}

\begin{quote}
\textbf{\sk{कियत्} \sk{दूरम्} \sk{अस्ति}?}
\end{quote}

\begin{quote}
\emph{kiyat dūram asti?} — \textquotedblleft{}How far is it?\textquotedblright{}
\end{quote}

\textbf{\sk{कियत्}} (\emph{kiyat}), \textquotedblleft{}how much\textquotedblright{}, with \textbf{\sk{दूरम्}} (\emph{dūram}), \textquotedblleft{}far\textquotedblright{}, asks the distance, and \textbf{\sk{कथं} \sk{गच्छामि}?} (\emph{kathaṃ gacchāmi?}) asks \textquotedblleft{}how do I get there?\textquotedblright{} The answer:

\begin{quote}
\textbf{\sk{बहु} \sk{दूरं} \sk{न} \sk{अस्ति।}}
\end{quote}

\begin{quote}
\emph{bahu dūraṃ na asti.} — \textquotedblleft{}It is not very far.\textquotedblright{}
\end{quote}

\subsection*{You'll want to know: left, right, straight on}
\textbf{\sk{वामतः}} (\emph{vāmataḥ}), \textquotedblleft{}on the left\textquotedblright{}, is \textbf{\sk{वामः}}, \textquotedblleft{}left\textquotedblright{}, with \emph{-taḥ}; \textbf{\sk{दक्षिणतः}} (\emph{dakṣiṇataḥ}), \textquotedblleft{}on the right\textquotedblright{}, is \textbf{\sk{दक्षिणः}} with it. \textbf{\sk{ऋजु}} (\emph{ṛju}), \textquotedblleft{}straight\textquotedblright{}, is \textbf{\sk{ऋजुः}} said as an adverb. What to do takes \textbf{-\sk{तु}}, the polite ending of \textbf{\sk{ददातु}} (\emph{dadātu}): \textbf{\sk{गच्छतु}} (\emph{gacchatu}), \textquotedblleft{}please go\textquotedblright{}.

\begin{quote}
\textbf{\sk{ऋजु} \sk{गच्छतु।}}
\end{quote}

\begin{quote}
\emph{ṛju gacchatu.} — \textquotedblleft{}Go straight on.\textquotedblright{}
\end{quote}

\begin{quote}
\textbf{\sk{ततः} \sk{वामतः} \sk{गच्छतु।}}
\end{quote}

\begin{quote}
\emph{tataḥ vāmataḥ gacchatu.} — \textquotedblleft{}Then turn left.\textquotedblright{}
\end{quote}

\begin{quote}
\textbf{\sk{स्थानकं} \sk{दक्षिणतः} \sk{अस्ति।}}
\end{quote}

\begin{quote}
\emph{sthānakaṃ dakṣiṇataḥ asti.} — \textquotedblleft{}The station is on the right.\textquotedblright{}
\end{quote}

\textbf{-\sk{तु}} speaks to \textbf{\sk{भवान्}}, or to \textbf{\sk{भवती}} (\emph{bhavatī}), a woman:

\begin{quote}
\textbf{\sk{भवती} \sk{अत्र} \sk{तिष्ठतु।}}
\end{quote}

\begin{quote}
\emph{bhavatī atra tiṣṭhatu.} — \textquotedblleft{}Please wait here.\textquotedblright{}
\end{quote}

\subsection*{Your turn}
Say these aloud:

\begin{itemize}
\raggedright
  \item \textquotedblleft{}Please go left\textquotedblright{} — \emph{kṛpayā vāmataḥ gacchatu}
  \item \textquotedblleft{}Please go right\textquotedblright{} — \emph{kṛpayā dakṣiṇataḥ gacchatu}
  \item \textquotedblleft{}Please go straight on\textquotedblright{} — \emph{kṛpayā ṛju gacchatu}
\end{itemize}

\begin{tcolorbox}[breakable,colback=teal!4,colframe=teal!35!black,title={Before you move on}]
\textquotedblleft{}Where is the station?\textquotedblright{} (\textbf{\sk{स्थानकं} \sk{कुत्र} \sk{अस्ति}?}) \textquotedblleft{}How far is it?\textquotedblright{} (\textbf{\sk{कियत्} \sk{दूरम्} \sk{अस्ति}?}) \textquotedblleft{}On the left\textquotedblright{}? (\textbf{\sk{वामतः}}.) \textquotedblleft{}Please go\textquotedblright{}? (\textbf{\sk{गच्छतु}}.)
\end{tcolorbox}
\section[(review)]{Review — asking the way}

\label{lesson:SA-R280-asking-the-way}



\begin{quote}
Say \textquotedblleft{}you should read every day\textquotedblright{}, to a friend. (\emph{tvaṃ pratidinaṃ paṭheḥ}.) Then \textquotedblleft{}therefore she is staying at home\textquotedblright{}. (\emph{tasmāt sā gṛhe tiṣṭhati}.)
\end{quote}

\subsection*{Your turn}
Say these aloud:

\begin{itemize}
\raggedright
  \item Please, where is the shop?
  \item How do I get to the school?
  \item How far is the station? It is quite near.
  \item Go straight on, then turn left.
  \item The station is on the right. Please wait here.
\end{itemize}

\begin{tcolorbox}[breakable,colback=teal!4,colframe=teal!35!black,title={Before you move on}]
\textbf{\sk{वामतः}} or \textbf{\sk{दक्षिणतः}}, for \textquotedblleft{}on the right\textquotedblright{}? (\textbf{\sk{दक्षिणतः}}.) Which ending makes a request polite? (\textbf{-\sk{तु}}: \textbf{\sk{गच्छतु}}, \textbf{\sk{तिष्ठतु}}.) Now the other way round: what does \textbf{\sk{कथं} \sk{गच्छामि}?} ask? (\textquotedblleft{}How do I get there?\textquotedblright{})
\end{tcolorbox}
